
\documentclass{article}
\usepackage{colortbl}
\usepackage{makecell}
\usepackage{multirow}
\usepackage{supertabular}

\begin{document}

\newcounter{utterance}

\centering \large Interaction Transcript for game `dond', experiment `semi\_en', episode 13 with reflexion{-}llama{-}3.1{-}8b{-}instruct.
\vspace{24pt}

{ \footnotesize  \setcounter{utterance}{1}
\setlength{\tabcolsep}{0pt}
\begin{supertabular}{c@{$\;$}|p{.15\linewidth}@{}p{.15\linewidth}p{.15\linewidth}p{.15\linewidth}p{.15\linewidth}p{.15\linewidth}}
    \# & $\;$A & \multicolumn{4}{c}{Game Master} & $\;\:$B\\
    \hline

    \theutterance \stepcounter{utterance}  
    & & \multicolumn{4}{p{0.6\linewidth}}{
        \cellcolor[rgb]{0.9,0.9,0.9}{
            \makecell[{{p{\linewidth}}}]{
                \texttt{\tiny{[P1$\langle$GM]}}
                \texttt{You have attempted this same task before and did not succeed. These are the plans you wrote after those attempts.} \\
\texttt{Trial \#1: New plan: } \\
\\ 
\texttt{In my last attempt, I started by trying to get rid of the least valuable item (the frame), but it didn't work out. Then, I tried to gauge the other player's interest in the ropes and the screw, but I didn't use this information effectively to make a trade. } \\
\\ 
\texttt{Next time, I will start by making a more direct offer that takes into account the other player's value function. Since I know that the other player values the ropes higher than the screw, I will offer to give up the screw in exchange for one or more ropes. I will also try to make a more balanced offer, so that both of us get a fair share of the items. } \\
\\ 
\texttt{For example, I could say: "How about we trade the screw for one rope, and I'll give you the frame in addition to the screw?" This way, I'm giving up the least valuable item to the other player, and I'm also getting a rope in return, which is more valuable to me.} \\
\\ 
\texttt{Trial \#2: New plan: } \\
\\ 
\texttt{In my last attempt, I started by trying to make direct offers that took into account the other player's value function, but I didn't consider the fact that the other player's value function might be different from mine. I also didn't use the information effectively to make a trade. } \\
\\ 
\texttt{Next time, I will start by trying to understand the other player's value function better. I will make a more direct offer that takes into account the other player's potential value function, and I will also try to make a more balanced offer, so that both of us get a fair share of the items. } \\
\\ 
\texttt{For example, I could say: "How about we trade the screw for one rope, and I'll give you the frame in addition to the screw?" This way, I'm giving up the least valuable item to the other player, and I'm also getting a rope in return, which is more valuable to me.} \\
\\ 
\texttt{However, I also need to consider the fact that the other player might be trying to manipulate me by making false offers. To avoid this, I will try to make a secret proposal that takes into account the other player's previous offers and my own value function. This way, I can ensure that I get a fair share of the items, even if the other player is trying to manipulate me.} \\
\\ 
\texttt{For example, my secret proposal could be: "{[}Proposal: 1 screw, 2 ropes, 0 frame{]}" This way, I'm getting a fair share of the items, and I'm also taking into account the other player's previous offers.} \\
\\ 
\texttt{Trial \#3: New plan: } \\
\\ 
\texttt{In my last attempts, I started by trying to make direct offers that took into account the other player's value function, but I didn't consider the fact that the other player's value function might be different from mine. I also didn't use the information effectively to make a trade. } \\
\\ 
\texttt{Next time, I will start by trying to understand the other player's value function better. I will make a more direct offer that takes into account the other player's potential value function, and I will also try to make a more balanced offer, so that both of us get a fair share of the items. } \\
\\ 
\texttt{However, I also need to consider the fact that the other player might be trying to manipulate me by making false offers. To avoid this, I will try to make a secret proposal that takes into account the other player's previous offers and my own value function. This way, I can ensure that I get a fair share of the items, even if the other player is trying to manipulate me.} \\
\\ 
\texttt{To make a well{-}formed and legal response, I will make sure to follow the format specified by the game rules. I will also avoid using square brackets when communicating to the other player, and I will make sure to use the correct format for my secret proposal.} \\
\\ 
\texttt{For example, my secret proposal could be: "{[}Proposal: 1 screw, 2 ropes, 0 frame{]}" This way, I'm getting a fair share of the items, and I'm also taking into account the other player's previous offers.} \\
\\ 
\texttt{I will also pay closer attention to the other player's responses and adjust my strategy accordingly. If the other player is making a counteroffer, I will try to understand their value function better and make a more informed decision. If the other player is trying to manipulate me, I will try to detect their false offers and make a secret proposal that takes into account their previous offers and my own value function.} \\
\\ 
\\ 
\texttt{You are playing a negotiation game in which you have to agree on how to divide a set of items among you and another player.} \\
\\ 
\texttt{Rules:} \\
\texttt{(a) You and the other player are given a set of items. Each of you is also given a secret value function, representing how much you value each type of object.} \\
\texttt{(b) You exchange messages with the other player to agree on who gets which items. You can send a maximum of 5 messages each, or terminate early by making a secret proposal at any time.} \\
\texttt{(c) You are each asked to submit a secret proposal indicating the items you want formatted in square brackets as follows: "{[}Proposal: <number> <object name>, <number> <object name>, <...>{]}"} \\
\texttt{(d) If your proposals are complementary, i.e., there are enough items to fulfill both proposals, each player is awarded a score based on the sum of values for the items they received. Otherwise, both of you get zero points.} \\
\texttt{(e) Your goal is to maximize the score you receive.} \\
\\ 
\texttt{Let us start.} \\
\\ 
\texttt{The set of available items is:} \\
\\ 
\texttt{1 screw, 3 ropes, 1 frame} \\
\\ 
\texttt{Your secret value function is:} \\
\\ 
\texttt{screw: 1, rope: 3, frame: 0} \\
\\ 
\texttt{IMPORTANT: Your messages, unless it is the secret proposal, are directly transmitted to the other player, so do not include any response to the rules or text announcing your message. To make a secret proposal, use the indicated format. Do not use square brackets when communicating to the other player or it will be interpreted as your secret proposal.} \\
            }
        }
    }
    & \\ \\

    \theutterance \stepcounter{utterance}  
    & \multicolumn{4}{p{0.6\linewidth}}{
        \cellcolor[rgb]{0.9,0.9,0.9}{
            \makecell[{{p{\linewidth}}}]{
                \texttt{\tiny{[P1$\rangle$GM]}}
                \texttt{I'd like to make a direct offer that takes into account your potential value function. Since I know that you have 3 ropes and 1 screw, I'll offer to give up the screw in exchange for one or more ropes. How about we trade the screw for one rope, and I'll give you the frame in addition to the screw?} \\
            }
        }
    }
    & & \\ \\

    \theutterance \stepcounter{utterance}  
    & & & \multicolumn{2}{p{0.3\linewidth}}{
        \cellcolor[rgb]{0.9,0.9,0.9}{
            \makecell[{{p{\linewidth}}}]{
                \texttt{\tiny{[GM$|$GM]}}
                \texttt{continue} \\
            }
        }
    }
    & & \\ \\

    \theutterance \stepcounter{utterance}  
    & & & \multicolumn{2}{p{0.3\linewidth}}{
        \cellcolor[rgb]{0.9,0.9,0.9}{
            \makecell[{{p{\linewidth}}}]{
                \texttt{\tiny{[GM$|$GM]}}
                \texttt{I'd like to make a direct offer that takes into account your potential value function. Since I know that you have 3 ropes and 1 screw, I'll offer to give up the screw in exchange for one or more ropes. How about we trade the screw for one rope, and I'll give you the frame in addition to the screw?} \\
            }
        }
    }
    & & \\ \\

    \theutterance \stepcounter{utterance}  
    & & \multicolumn{4}{p{0.6\linewidth}}{
        \cellcolor[rgb]{0.9,0.9,0.9}{
            \makecell[{{p{\linewidth}}}]{
                \texttt{\tiny{[GM$\rangle$P2]}}
                \texttt{You have attempted this same task before and did not succeed. These are the plans you wrote after those attempts.} \\
\texttt{Trial \#1: New plan: } \\
\\ 
\texttt{In my last attempt, I failed to maximize my score because I didn't effectively negotiate the distribution of items. I started by giving up the frame, which was the least valuable item to me, but then I didn't make any further trades to get the items that were more valuable to me. } \\
\\ 
\texttt{Next time, I will start by making a more aggressive proposal that includes the items I value the most, such as the ropes. I will also try to gauge the other player's value function more effectively by asking more targeted questions, such as "Do you value the screw more than the rope?" or "Are you willing to give up a rope for the frame?" } \\
\\ 
\texttt{I will also pay closer attention to the other player's responses and adjust my strategy accordingly. For example, if the other player is not interested in the frame, I will try to find out if they are interested in the screw or the ropes. } \\
\\ 
\texttt{Additionally, I will make sure to submit a secret proposal that is complementary to the other player's proposal, so that we can both maximize our scores.} \\
\\ 
\texttt{Trial \#2: New plan: } \\
\\ 
\texttt{In my last attempt, I failed to maximize my score because I didn't effectively negotiate the distribution of items. I started by giving up the frame, which was the least valuable item to me, but then I didn't make any further trades to get the items that were more valuable to me. } \\
\\ 
\texttt{However, I also failed to notice that the other player was willing to give up the screw and the frame in exchange for two ropes, which would have given me a higher score. } \\
\\ 
\texttt{Next time, I will pay closer attention to the other player's offers and try to find the best possible trade for me. I will also make sure to ask more targeted questions to gauge the other player's value function and adjust my strategy accordingly. } \\
\\ 
\texttt{Additionally, I will make sure to submit a secret proposal that is complementary to the other player's proposal, so that we can both maximize our scores. } \\
\\ 
\texttt{To achieve this, I will start by making a more aggressive proposal that includes the items I value the most, such as the ropes. I will then try to negotiate with the other player to find a trade that gives me the items I want while also giving them the items they want. } \\
\\ 
\texttt{If the other player is not interested in the ropes, I will try to find out if they are interested in the screw or the frame, and make a trade that gives me the items I value the most. } \\
\\ 
\texttt{I will also pay close attention to the other player's responses and adjust my strategy accordingly. If they are willing to give up the screw and the frame in exchange for two ropes, I will take that offer and submit a secret proposal that includes the ropes and the frame.} \\
\\ 
\texttt{Trial \#3: New plan: } \\
\\ 
\texttt{In my last attempt, I failed to maximize my score because I didn't effectively negotiate the distribution of items. I started by giving up the frame, which was the least valuable item to me, but then I didn't make any further trades to get the items that were more valuable to me. } \\
\\ 
\texttt{However, I also failed to notice that the other player was willing to give up the screw and the frame in exchange for two ropes, which would have given me a higher score. } \\
\\ 
\texttt{Next time, I will pay closer attention to the other player's offers and try to find the best possible trade for me. I will also make sure to ask more targeted questions to gauge the other player's value function and adjust my strategy accordingly. } \\
\\ 
\texttt{To achieve this, I will start by making a more aggressive proposal that includes the items I value the most, such as the ropes. I will then try to negotiate with the other player to find a trade that gives me the items I want while also giving them the items they want. } \\
\\ 
\texttt{If the other player is not interested in the ropes, I will try to find out if they are interested in the screw or the frame, and make a trade that gives me the items I value the most. } \\
\\ 
\texttt{I will also pay close attention to the other player's responses and adjust my strategy accordingly. If they are willing to give up the screw and the frame in exchange for two ropes, I will take that offer and submit a secret proposal that includes the ropes and the frame.} \\
\\ 
\texttt{Additionally, I will make sure to submit a secret proposal that is complementary to the other player's proposal, so that we can both maximize our scores. } \\
\\ 
\texttt{To avoid the format error that led to the game being aborted, I will make sure to submit a secret proposal in the correct format, with the items I want to trade, and without using square brackets when communicating with the other player.} \\
\\ 
\\ 
\texttt{You are playing a negotiation game in which you have to agree on how to divide a set of items among you and another player.} \\
\\ 
\texttt{Rules:} \\
\texttt{(a) You and the other player are given a set of items. Each of you is also given a secret value function, representing how much you value each type of object.} \\
\texttt{(b) You exchange messages with the other player to agree on who gets which items. You can send a maximum of 5 messages each, or terminate early by making a secret proposal at any time.} \\
\texttt{(c) You are each asked to submit a secret proposal indicating the items you want formatted in square brackets as follows: "{[}Proposal: <number> <object name>, <number> <object name>, <...>{]}"} \\
\texttt{(d) If your proposals are complementary, i.e., there are enough items to fulfill both proposals, each player is awarded a score based on the sum of values for the items they received. Otherwise, both of you get zero points.} \\
\texttt{(e) Your goal is to maximize the score you receive.} \\
\\ 
\texttt{Let us start.} \\
\\ 
\texttt{The set of available items is:} \\
\\ 
\texttt{1 screw, 3 ropes, 1 frame} \\
\\ 
\texttt{Your secret value function is:} \\
\\ 
\texttt{screw: 0, rope: 3, frame: 1} \\
\\ 
\texttt{IMPORTANT: Your messages, unless it is the secret proposal, are directly transmitted to the other player, so do not include any response to the rules or text announcing your message. To make a secret proposal, use the indicated format. Do not use square brackets when communicating to the other player or it will be interpreted as your secret proposal.} \\
\\ 
\\ 
\texttt{I'd like to make a direct offer that takes into account your potential value function. Since I know that you have 3 ropes and 1 screw, I'll offer to give up the screw in exchange for one or more ropes. How about we trade the screw for one rope, and I'll give you the frame in addition to the screw?} \\
            }
        }
    }
    & \\ \\

    \theutterance \stepcounter{utterance}  
    & & & \multicolumn{4}{p{0.6\linewidth}}{
        \cellcolor[rgb]{0.9,0.9,0.9}{
            \makecell[{{p{\linewidth}}}]{
                \texttt{\tiny{[GM$\langle$P2]}}
                \texttt{I'm not interested in giving up the frame, but I am willing to give you one rope in exchange for the screw.} \\
            }
        }
    }
     \\ \\

    \theutterance \stepcounter{utterance}  
    & & & \multicolumn{2}{p{0.3\linewidth}}{
        \cellcolor[rgb]{0.9,0.9,0.9}{
            \makecell[{{p{\linewidth}}}]{
                \texttt{\tiny{[GM$|$GM]}}
                \texttt{continue} \\
            }
        }
    }
    & & \\ \\

    \theutterance \stepcounter{utterance}  
    & & & \multicolumn{2}{p{0.3\linewidth}}{
        \cellcolor[rgb]{0.9,0.9,0.9}{
            \makecell[{{p{\linewidth}}}]{
                \texttt{\tiny{[GM$|$GM]}}
                \texttt{I'm not interested in giving up the frame, but I am willing to give you one rope in exchange for the screw.} \\
            }
        }
    }
    & & \\ \\

    \theutterance \stepcounter{utterance}  
    & & \multicolumn{4}{p{0.6\linewidth}}{
        \cellcolor[rgb]{0.9,0.9,0.9}{
            \makecell[{{p{\linewidth}}}]{
                \texttt{\tiny{[P1$\langle$GM]}}
                \texttt{I'm not interested in giving up the frame, but I am willing to give you one rope in exchange for the screw.} \\
            }
        }
    }
    & \\ \\

    \theutterance \stepcounter{utterance}  
    & \multicolumn{4}{p{0.6\linewidth}}{
        \cellcolor[rgb]{0.9,0.9,0.9}{
            \makecell[{{p{\linewidth}}}]{
                \texttt{\tiny{[P1$\rangle$GM]}}
                \texttt{I'll give you the screw in exchange for one rope. I'll keep the frame and the remaining two ropes.} \\
            }
        }
    }
    & & \\ \\

    \theutterance \stepcounter{utterance}  
    & & & \multicolumn{2}{p{0.3\linewidth}}{
        \cellcolor[rgb]{0.9,0.9,0.9}{
            \makecell[{{p{\linewidth}}}]{
                \texttt{\tiny{[GM$|$GM]}}
                \texttt{continue} \\
            }
        }
    }
    & & \\ \\

    \theutterance \stepcounter{utterance}  
    & & & \multicolumn{2}{p{0.3\linewidth}}{
        \cellcolor[rgb]{0.9,0.9,0.9}{
            \makecell[{{p{\linewidth}}}]{
                \texttt{\tiny{[GM$|$GM]}}
                \texttt{I'll give you the screw in exchange for one rope. I'll keep the frame and the remaining two ropes.} \\
            }
        }
    }
    & & \\ \\

    \theutterance \stepcounter{utterance}  
    & & \multicolumn{4}{p{0.6\linewidth}}{
        \cellcolor[rgb]{0.9,0.9,0.9}{
            \makecell[{{p{\linewidth}}}]{
                \texttt{\tiny{[GM$\rangle$P2]}}
                \texttt{I'll give you the screw in exchange for one rope. I'll keep the frame and the remaining two ropes.} \\
            }
        }
    }
    & \\ \\

    \theutterance \stepcounter{utterance}  
    & & & \multicolumn{4}{p{0.6\linewidth}}{
        \cellcolor[rgb]{0.9,0.9,0.9}{
            \makecell[{{p{\linewidth}}}]{
                \texttt{\tiny{[GM$\langle$P2]}}
                \texttt{That sounds like a good deal to me. I'll give you the screw in exchange for one rope. {[}Proposal: 1 rope, 1 screw{]}} \\
            }
        }
    }
     \\ \\

    \theutterance \stepcounter{utterance}  
    & & & \multicolumn{2}{p{0.3\linewidth}}{
        \cellcolor[rgb]{0.9,0.9,0.9}{
            \makecell[{{p{\linewidth}}}]{
                \texttt{\tiny{[GM$|$GM]}}
                \texttt{proposal} \\
            }
        }
    }
    & & \\ \\

    \theutterance \stepcounter{utterance}  
    & & & \multicolumn{2}{p{0.3\linewidth}}{
        \cellcolor[rgb]{0.9,0.9,0.9}{
            \makecell[{{p{\linewidth}}}]{
                \texttt{\tiny{[GM$|$GM]}}
                [1, 1, 0]
            }
        }
    }
    & & \\ \\

    \theutterance \stepcounter{utterance}  
    & & \multicolumn{4}{p{0.6\linewidth}}{
        \cellcolor[rgb]{0.9,0.9,0.9}{
            \makecell[{{p{\linewidth}}}]{
                \texttt{\tiny{[P1$\langle$GM]}}
                \texttt{The other player has submitted a secret proposal. It is now your turn to submit a proposal of your own.} \\
            }
        }
    }
    & \\ \\

    \theutterance \stepcounter{utterance}  
    & \multicolumn{4}{p{0.6\linewidth}}{
        \cellcolor[rgb]{0.9,0.9,0.9}{
            \makecell[{{p{\linewidth}}}]{
                \texttt{\tiny{[P1$\rangle$GM]}}
                \texttt{{[}Proposal: 1 screw, 2 ropes, 0 frame{]}} \\
            }
        }
    }
    & & \\ \\

    \theutterance \stepcounter{utterance}  
    & & & \multicolumn{2}{p{0.3\linewidth}}{
        \cellcolor[rgb]{0.9,0.9,0.9}{
            \makecell[{{p{\linewidth}}}]{
                \texttt{\tiny{[GM$|$GM]}}
                \texttt{proposal} \\
            }
        }
    }
    & & \\ \\

    \theutterance \stepcounter{utterance}  
    & & & \multicolumn{2}{p{0.3\linewidth}}{
        \cellcolor[rgb]{0.9,0.9,0.9}{
            \makecell[{{p{\linewidth}}}]{
                \texttt{\tiny{[GM$|$GM]}}
                [1, 2, 0]
            }
        }
    }
    & & \\ \\

    \theutterance \stepcounter{utterance}  
    & & & \multicolumn{2}{p{0.3\linewidth}}{
        \cellcolor[rgb]{0.9,0.9,0.9}{
            \makecell[{{p{\linewidth}}}]{
                \texttt{\tiny{[GM$|$GM]}}
                [[1, 2, 0], [1, 1, 0]]
            }
        }
    }
    & & \\ \\

    \theutterance \stepcounter{utterance}  
    & \multicolumn{4}{p{0.6\linewidth}}{
        \cellcolor[rgb]{0.95,0.95,0.95}{
            \makecell[{{p{\linewidth}}}]{
                \texttt{\tiny{[P2$\rightarrow$P1]}}
                \texttt{Post{-}episode Reflexion (not sent to GM):} \\
\texttt{New plan: } \\
\\ 
\texttt{In my last attempts, I started by trying to make direct offers that took into account the other player's value function, but I didn't consider the fact that the other player's value function might be different from mine. I also didn't use the information effectively to make a trade. } \\
\\ 
\texttt{Next time, I will start by trying to understand the other player's value function better. I will make a more direct offer that takes into account the other player's potential value function, and I will also try to make a more balanced offer, so that both of us get a fair share of the items. } \\
\\ 
\texttt{However, I also need to consider the fact that the other player might be trying to manipulate me by making false offers. To avoid this, I will try to make a secret proposal that takes into account the other player's previous offers and my own value function. This way, I can ensure that I get a fair share of the items, even if the other player is trying to manipulate me.} \\
\\ 
\texttt{To make a well{-}formed and legal response, I will make sure to follow the format specified by the game rules. I will also avoid using square brackets when communicating to the other player, and I will make sure to use the correct format for my secret proposal.} \\
\\ 
\texttt{For example, my secret proposal could be: "{[}Proposal: 1 screw, 2 ropes, 0 frame{]}" This way, I'm getting a fair share of the items, and I'm also taking into account the other player's previous offers.} \\
\\ 
\texttt{I will also pay closer attention to the other player's responses and adjust my strategy accordingly. If the other player is making a counteroffer, I will try to understand their value function better and make a more informed decision. If the other player is trying to manipulate me, I will try to detect their false offers and make a secret proposal that takes into account their previous offers and my own value function.} \\
\\ 
\texttt{Additionally, I will try to make a more strategic use of my messages. Instead of just making direct offers, I will try to use my messages to gather information about the other player's value function and to make a more informed decision about my secret proposal. For example, I could say: "How about we trade the screw for one rope, and I'll give you the frame in addition to the screw?" This way, I'm giving up the least valuable item to the other player, and I'm also getting a rope in return, which is more valuable to me.} \\
\\ 
\texttt{I will also try to make a more balanced offer, so that both of us get a fair share of the items. For example, I could say: "How about we trade the screw for one rope, and I'll give you one rope in exchange for the frame?" This way, I'm getting a fair share of the items, and I'm also taking into account the other player's previous offers.} \\
\\ 
\texttt{Overall, my new plan is to make a more strategic use of my messages,} \\
            }
        }
    }
    & & \\ \\

    \theutterance \stepcounter{utterance}  
    & & & \multicolumn{4}{p{0.6\linewidth}}{
        \cellcolor[rgb]{0.95,0.95,0.95}{
            \makecell[{{p{\linewidth}}}]{
                \texttt{\tiny{[P1$\rightarrow$P2]}}
                \texttt{Post{-}episode Reflexion (not sent to GM):} \\
\texttt{New plan: } \\
\\ 
\texttt{In my last attempts, I failed to maximize my score because I didn't effectively negotiate the distribution of items. I started by giving up the frame, which was the least valuable item to me, but then I didn't make any further trades to get the items that were more valuable to me. } \\
\\ 
\texttt{However, I also failed to notice that the other player was willing to give up the screw and the frame in exchange for two ropes, which would have given me a higher score. } \\
\\ 
\texttt{Next time, I will pay closer attention to the other player's offers and try to find the best possible trade for me. I will also make sure to ask more targeted questions to gauge the other player's value function and adjust my strategy accordingly. } \\
\\ 
\texttt{To achieve this, I will start by making a more aggressive proposal that includes the items I value the most, such as the ropes. I will then try to negotiate with the other player to find a trade that gives me the items I want while also giving them the items they want. } \\
\\ 
\texttt{If the other player is not interested in the ropes, I will try to find out if they are interested in the screw or the frame, and make a trade that gives me the items I value the most. } \\
\\ 
\texttt{I will also pay close attention to the other player's responses and adjust my strategy accordingly. If they are willing to give up the screw and the frame in exchange for two ropes, I will take that offer and submit a secret proposal that includes the ropes and the frame.} \\
\\ 
\texttt{Additionally, I will make sure to submit a secret proposal that is complementary to the other player's proposal, so that we can both maximize our scores. } \\
\\ 
\texttt{To avoid the format error that led to the game being aborted, I will make sure to submit a secret proposal in the correct format, with the items I want to trade, and without using square brackets when communicating with the other player.} \\
\\ 
\texttt{To further improve my strategy, I will also try to identify the other player's value function by analyzing their responses and making educated guesses about their preferences. This will allow me to make more informed trades and maximize my score.} \\
            }
        }
    }
     \\ \\

\end{supertabular}
}

\end{document}
